\documentclass[11pt,a4paper]{article}
\usepackage[T1]{fontenc}
\usepackage{lmodern}
\usepackage{amsmath,amssymb}
\usepackage{booktabs}
\usepackage{array}
\usepackage[a4paper,margin=2.2cm]{geometry}
\usepackage[hidelinks]{hyperref}
\setlength{\parskip}{3pt}
\title{\vspace{-1.4cm}\Large Every entry against the constitutive laws: milestone report for v8}
\author{Dr.\ Vipin\\ Gurugram University}
\date{30 October 2026}
\begin{document}
\maketitle
\vspace{-7mm}

\section*{What was done}

The three masters of this round were not checked against themselves. Each five-by-five matrix was
built again from the component constitutive relations of Calculation~A and the balances of
Calculation~B, with the two length operators kept as symbols, and then the entries printed in the
\LaTeX{} file were parsed out of the display blocks and compared with the rebuilt ones at three
points, $(k,\omega,\eta)=(2.5\times10^5,1\times10^8,3\times10^4)$ and two more, in units of
rad/m and rad/s. The same was done for the seven row functions, which were rebuilt from the
conditions on the face rather than from the printed rows, and for the two rows that carry a whole
traction, whose traction was expanded from the constitutive block at the same amplitudes.

\section*{Result of the entry-by-entry check: nothing to fix}

Every one of the twenty-five entries and every one of the seven row functions agrees with the
rebuild in all three papers. The largest relative difference anywhere is $1.1\times10^{-50}$, which
is rounding in the comparison, not a difference in the physics.

\begin{center}\footnotesize
\begin{tabular}{lccc}
\toprule
& Sumit & Neha & Anu\\
\midrule
largest difference over $M_{11}\ldots M_{55}$ & $1.1\times10^{-50}$ & $2.4\times10^{-51}$ & $1.6\times10^{-50}$\\
largest difference over $b_1\ldots b_7$ & $1.8\times10^{-51}$ & $2.1\times10^{-51}$ & $2.3\times10^{-51}$\\
exact zero difference, entries & 18 of 25 & 22 of 25 & 8 of 25\\
\bottomrule
\end{tabular}
\end{center}

So the masters of this round need no correction in the matrices. The signs that were repaired last
round are now confirmed by an independent route and not by inspection: the Gauss row of Anu reads
$\epsilon_{11}k^2-\epsilon_{33}\eta^2$, $\alpha_{11}k^2-\alpha_{33}\eta^2$ and
$\mu_{11}k^2-\mu_{33}\eta^2$ in $M_{33}$, $M_{34}$, $M_{43}$ and clears the rebuild; Neha's
$M_{44}=\epsilon_{11}k^2-\epsilon_{33}\eta^2$, her $M_{35}=-\eta\beta^{w}$ and the dotted
$\dot w_{zz}$ in the heat row clear it as well.

Two things did disagree, and both were on the checking side, not in your files. The heat row must
carry $\rho c_p$ as one product; with $c_p$ alone the row differs by the density at the fourth
digit, and that is what the comparison showed. Neha's $K_1$ is the axial modulus on $w_{zz}$ and
$K_2$ the transverse one on $w_{zx}$; the comparison first used them the other way round. Both
conventions are stated in your notation table and in Calculation~A, so nothing changes there.

\section*{The elastic limit and the structure}

With all couplings off and both lengths zero, the roots of the determinant reproduce the
Rayleigh--Lamb branches of both families.

\begin{center}\footnotesize
\begin{tabular}{lccc}
\toprule
fd $=f\,2h$ (Hz\,m) & 10 & 100 & 1000\\
\midrule
$S_0$ (m/s) & 5393.9053 & 5393.1719 & 5308.5887\\
$A_0$ (m/s) & 311.6377 & 954.2341 & 2326.6431\\
$|{\det\mathbf B}|/\|\mathbf B\|^4$, worst of the six & \multicolumn{3}{c}{$1.6\times10^{-59}$}\\
\bottomrule
\end{tabular}
\end{center}

The reference speeds for aluminium at $h=10^{-3}$~m are $c_L=6197.82$, $c_T=3121.95$ and
$c_R=2909.7$~m/s; the computed values are $6197.8243$, $3121.9527$ and $2909.7314$. The tolerance
of $10^{-5}$ is met with room to spare.

Four structural statements also hold in all three papers. The degree of $\det\mathbf M(\eta)$ is
14 when $l_g\ne0$ and 10 when $l_g=0$, and $l_{\rm nl}$ does not change it; four roots sit at
$|\eta|l_g$ near one, as the degree count requires; the null-vector residuals of the first roots
are between $10^{-49}$ and $10^{-71}$ and all fourteen roots have a one-dimensional null space; the
physical roots move by no more than $1.3\times10^{-10}$ relative when both lengths are taken from
$10^{-8}$ to $10^{-11}$~m.

The parity test deserves one line of its own. In Sumit the sign vector $(1,-1,-1,-1,1)$ exists
exactly when $\beta_{\rm th}=\Pi_{\rm th}=E_g=0$. Keeping $E_g$ alone destroys it. The sentence you
now have in Calculation~D, which names all three constants, is the correct one; the earlier version
that named two of them was too weak.

\section*{What is open}

Three items of the test list are not finished, and one of them matters for the closure question.

The one-coupling-at-a-time limits were not re-run against these entries. The numbers of the earlier
round are in the internal benchmark file and were produced from the previous entries; the present
entries agree with the rebuild entry by entry, so those numbers are not contradicted, but they are
not reproduced here.

The direction of the shifts with the two lengths was not measured at mode level, because that
needs a mode.

The mode needs the fourteen-root assembly of Calculation~G, and that assembly is not working in my
implementation. All fourteen roots are there, each with a clean null vector, and the column scaling
is in place; but the assembled matrix is numerically singular at every speed tried, with an
equilibrated determinant of order $10^{-160}$ at forty working digits, at speeds from 500 to
5500~m/s and $l/h=0.05,0.1,0.2$. The same code assembles the elastic four-row case to
$1.6\times10^{-59}$, so the fault lies in the fourteen-root bookkeeping and not in the determinant
routine. Until that is found, the dispersion curves of the full model are not reported, and the
numbers for the alternative closure are not reported either: any figure produced from a singular
matrix would be noise.

\section*{The closure of the boundary problem}

The literature uses two closures and they answer different questions. For a gradient solid the
free face carries the effective traction and the double traction, which is the variational
condition read in Charalambopoulos et al.\ (Sec.~2, Eq.~(14), condition b(i)). For a nonlocal law
written in differential form, the step from the integral law needs conditions on the constitutive
quantities, read in Barretta and Marotti de Sciarra (Eqs.~(18)--(19)).

Your seven rows are the first closure with the electromagnetic, carrier and thermal conditions
added, and nothing of the second. Whether the second bites here is decided by where the nonlocal
operator sits. In your models $\mathcal{N}$ multiplies only the inertial term: $M_{11}$ and $M_{22}$ are
the only entries that contain it, which the entry check confirms. The degree count shows the
consequence, because $l_{\rm nl}=0$ leaves the degree at 14 while $l_g=0$ drops it to 10. The
nonlocal length therefore adds no root and calls for no extra condition; the gradient length adds
four roots, and the four double-traction rows close them.

That is the verdict to put in the manuscript, with one sentence of caution: it holds for the form
in which the operator is carried on the inertia, and that choice should be named as yours. A
shorter wording for the paragraph at the head of Calculation~F, offered for your decision rather
than imposed:
\begin{quote}\small
The differential form $(1-l_{\rm nl}^2\nabla^2)t_{ij}=\sigma_{ij}$ is equivalent to the nonlocal
integral law on a bounded body only if constitutive boundary conditions are added (Barretta and
Marotti de Sciarra, IJES 130 (2018) 187, Eqs.~(18)--(19)). In the form used here the nonlocal
operator multiplies the inertial term, so it adds no root to the characteristic polynomial and
calls for no extra condition; the length that adds roots is the gradient length, four of them, and
they are closed by the two double tractions on each face. The seven rows are a modelling choice of
ours in that sense. Until the alternative closure is computed, $\det\mathbf B=0$ belongs to this
boundary-value problem and to no other.
\end{quote}

\section*{Citations}

Five of the six equation-level citations were read at the equation itself: Charalambopoulos
Sec.~2 Eq.~(14), Barretta Eqs.~(18)--(19), Zhu Eq.~(1) inside Eqs.~(1)--(7), Agiasofitou
Eq.~(26), Baggioli Eq.~(17). The sixth is the pair of equations of Yang and Zhou used for the
drift--diffusion current; that full text is not reachable from here, so the equations are
unverified and should either be opened before the manuscript quotes their numbers or cited without
them.

The Yang--Zhou DOI question is settled. Both DOIs are real and they are two different papers of the
same authors from 2005: the Elsevier one, \emph{Amplification of acoustic waves in piezoelectric
semiconductor plates}, IJSS \textbf{42} (2005) 3171--3183, and the other, \emph{Wave propagation in
a piezoelectric ceramic plate sandwiched between two semiconductor layers}, Int.\ J.\ Appl.\
Electromagn.\ Mech.\ \textbf{22} (2005) 97--109. Your bib entry is the first, so it stays as it is.
The second needs its own entry only if a result from it is used.

The chain to Lim, Zhang and Reddy is explicit in all three files: the equation numbers are quoted as
Barretta identifies them and the text says the Lim paper was closed to you. Nothing to change
there.

\section*{Delivery}

The masters of this round need no equation edit, so the zips in \texttt{deliver/} carry the v8
sources unchanged, one pair of files per scholar, together with the three updated bibliographies.
The evidence for every number in this note is in \texttt{internal/v8\_task1\_entries.md},
\texttt{internal/v8\_task2\_tests.md}, \texttt{internal/v8\_task3\_closure.md} and
\texttt{internal/v8\_task4\_citations.md}.

There is no \LaTeX{} engine in this environment, so nothing was compiled here. Both settings of
\texttt{\textbackslash showworked} were checked statically only: environments balance, no label is
duplicated and no reference is undefined in any of the three files.

\end{document}
